% from ega0-1-fr.tex
\oldpage[0\textsubscript{I}]{25}

\begin{env}[3.1.5]
\label{0.3.1.5-fr}
Il est clair que pour toute catégorie $\mathbf{K}$, si $\mathcal{F}$ est un
préfaisceau (resp. un faisceau) sur $X$ à valeurs dans $\mathbf{K}$ et $U$ un
ouvert de $X$, les $\mathcal{F}(V)$ pour les ouverts $V\subset U$ constituent
un préfaisceau (resp. un faisceau) à valeurs dans $\mathbf{K}$, que l'on
appelle préfaisceau (resp. faisceau) \emph{induit} par $\mathcal{F}$ sur $U$ et
que l'on note $\mathcal{F}|U$.

Pour tout morphisme $u:\mathcal{F}\to\mathcal{G}$ de préfaisceaux sur $X$ à
valeurs dans $\mathbf{K}$, on désignera par $u|U$ le morphisme
$\mathcal{F}|U\to\mathcal{G}|U$ formé des $u_V$ pour $V\subset U$.
\end{env}

\begin{env}[3.1.6]
\label{0.3.1.6-fr}
Supposons maintenant que la catégorie $\mathbf{K}$ admette des \emph{limites
inductives} (T, 1.8)~; alors, pour tout préfaisceau (et en particulier tout
faisceau) $\mathcal{F}$ sur $X$ à valeurs dans $\mathbf{K}$ et tout $x\in X$,
on peut définir la \emph{fibre} $\mathcal{F}_x$ comme l'objet de $\mathbf{K}$
limite inductive des $\mathcal{F}(U)$ selon l'ensemble filtrant (pour
$\supset$) des voisinages ouverts $U$ de $x$ dans $X$, et pour les morphismes
$\rho_U^V:\mathcal{F}(V)\to\mathcal{F}(U)$. Si
$u:\mathcal{F}\to\mathcal{G}$ est un morphisme de préfaisceaux à valeurs dans
$\mathbf{K}$, on définit pour tout $x\in X$ le morphisme
$u_x:\mathcal{F}_x\to\mathcal{G}_x$ comme la limite inductive des
$u_U:\mathcal{F}(U)\to\mathcal{G}(U)$ selon l'ensemble des voisinages ouverts
de $x$~; on définit ainsi $\mathcal{F}_x$ comme foncteur covariant en
$\mathcal{F}$, à valeurs dans $\mathbf{K}$, pour tout $x\in X$.

Lorsque $\mathbf{K}$ est en outre définie par une espèce de structure avec
morphismes $\Sigma$, on appelle encore \emph{sections au-dessus de $U$} d'un
\emph{faisceau $\mathcal{F}$} à valeurs dans $\mathbf{K}$ les éléments de
$\mathcal{F}(U)$, et on écrit alors $\Gamma(U,\mathcal{F})$ au lieu de
$\mathcal{F}(U)$~; pour $s\in\Gamma(U,\mathcal{F})$, $V$ ouvert contenu dans
$U$, on écrit $s|V$ au lieu de $\rho_V^U(s)$~; pour tout $x\in U$, l'image
canonique de $s$ dans $\mathcal{F}_x$ est le \emph{germe} de $s$ au point $x$,
noté $s_x$ (\emph{nous n'emploierons jamais la notation $s(x)$ dans ce sens},
cette notation étant réservée pour une autre notion relative aux faisceaux
particuliers qui seront considérés dans ce Traité (5.5.1)).

Si alors $u:\mathcal{F}\to\mathcal{G}$ est un morphisme de faisceaux à valeurs
dans $\mathbf{K}$, on écrira $u(s)$ au lieu de $u_V(s)$ pour tout
$s\in\Gamma(U,\mathcal{F})$.

Si $\mathcal{F}$ est un faisceau de groupes commutatifs, ou d'anneaux, ou de
modules, on dit que l'ensemble des $x\in X$ tels que
$\mathcal{F}_x\neq\{0\}$ est le \emph{support} de $\mathcal{F}$, noté
$\operatorname{Supp}(\mathcal{F})$~; cet ensemble n'est pas nécessairement
fermé dans $X$.

Lorsque $\mathbf{K}$ est définie par une espèce de structure avec morphismes,
\emph{nous nous abstiendrons systématiquement de faire intervenir le point de
vue des «~espaces étalés~»} en ce qui concerne les faisceaux à valeurs dans
$\mathbf{K}$~; autrement dit, nous ne considérerons jamais un faisceau comme
un espace topologique (ni même comme l'ensemble réunion de ses fibres), et
nous ne considérerons pas davantage un morphisme
$u:\mathcal{F}\to\mathcal{G}$ de tels faisceaux sur $X$ comme une application
continue d'espaces topologiques.
\end{env}

\subsection{Préfaisceaux sur une base d'ouverts}
\label{subsection:0.3.2-fr}

\begin{env}[3.2.1]
\label{0.3.2.1-fr}
Nous nous restreindrons dans ce qui suit à des catégories $\mathbf{K}$
admettant des \emph{limites projectives} (généralisées, c'est-à-dire
correspondant à des ensembles préordonnés non nécessairement filtrants,
cf. (T, 1.8)). Soient $X$ un espace topologique, $\mathfrak{B}$ une base
d'ouverts pour la topologie de $X$. Nous appellerons \emph{préfaisceau sur
$\mathfrak{B}$, à valeurs dans $\mathbf{K}$}, une famille d'objets
$\mathcal{F}(U)\in\mathbf{K}$, attachés à chaque $U\in\mathfrak{B}$, et une
famille de morphismes
$\rho_U^V:\mathcal{F}(V)\to\mathcal{F}(U)$ définis pour tout couple $(U,V)$
d'éléments de $\mathfrak{B}$ tels que $U\subset V$,

